\chapter[Równanie ciepła]{Równanie ciepła: jądro Gaussa, maksimum, wygładzanie i ewolucja}
\label{chap:pde-heat}

W poprzednim rozdziale warunek $\Delta u=0$ wymuszał silną
sztywność funkcji. Teraz badamy równanie, w którym laplasjan
steruje zmianą w czasie:
\begin{equation}\label{eq:pde4-heat}
 \partial_tu(t,x)-\kappa\Delta_xu(t,x)=f(t,x),\qquad
 t>0,\quad x\in\R^n,\quad \kappa>0.
\end{equation}
Gdy $f=0$, równanie opisuje swobodną dyfuzję. Fundamentalne pytania
są odmienne niż dla równania Laplace'a: jak odzyskać rozwiązanie
z profilu początkowego, jakie informacje zostają zachowane i dlaczego
każdy dodatni czas poprawia regularność? Wyprowadzimy rozwiązanie
całego zagadnienia Cauchy'ego, zasady maksimum, formułę Duhamela,
bilanse energii i przykłady pokazujące, że bez określenia klasy
rozwiązań nawet jednoznaczność może zawieść.

W całym rozdziale $n\geq1$ jest liczbą całkowitą,
$\kappa>0$, a funkcje skalarne są rzeczywiste, o ile
nie zaznaczono inaczej. Przez $\lambda_k$ oznaczamy
miarę Lebesgue'a w $\R^k$ dla $k\geq1$.
Dla $1\leq p<\infty$ symbole $L^p(\R^n)$ i $\|\cdot\|_p$
oznaczają przestrzeń klas funkcji mierzalnych
o całkowalnej $p$-tej potędze modułu oraz jej normę.
Przestrzeń $L^\infty(\R^n)$ składa się z klas funkcji
istotnie ograniczonych, a jej normą jest istotne
supremum modułu. Przez $\mathrm{BUC}(\R^n)$ oznaczamy
przestrzeń funkcji ograniczonych i jednostajnie
ciągłych z normą supremum. W dalszych twierdzeniach
\emph{silna ciągłość} oznacza zbieżność w normie danej
przestrzeni, a nie tylko punktową ani słabą. Będziemy korzystać z elementarnych
twierdzeń o całkowaniu: Fubiniego--Tonellego, zbieżności
zdominowanej, nierówności Minkowskiego i ciągłości przesunięć
w $L^p$; w rozdziałach o analizie funkcjonalnej wrócimy
do ich szerszego kontekstu.

\section{Jądro Gaussa: skąd bierze się rozwiązanie?}
\label{sec:pde4-gauss}

\begin{definicja}[Jądro cieplne]
\label{def:pde4-kernel}
Niech $n\geq1$ i $\kappa>0$. Dla $t>0$ oraz $x\in\R^n$
definiujemy
\begin{equation}\label{eq:pde4-kernel}
 \Gamma_\kappa(t,x)=\frac{1}{(4\pi\kappa t)^{n/2}}
 \exp\!\left(-\frac{|x|^2}{4\kappa t}\right).
\end{equation}
Zapis $\Gamma_\kappa(t)*g$ oznacza splot w zmiennej przestrzennej:
$(\Gamma_\kappa(t)*g)(x)=\int_{\R^n}\Gamma_\kappa(t,x-y)g(y)
\,\mathrm d\lambda_n(y)$, jeżeli całka ma sens.
\end{definicja}

\begin{twierdzenie}[Równanie, normalizacja i skalowanie]
\label{thm:pde4-gauss-basic}
Dla każdego $n\geq1$, $\kappa>0$, $t>0$ zachodzi:
$\Gamma_\kappa(t,x)>0$,
$(\partial_t-\kappa\Delta_x)\Gamma_\kappa(t,x)=0$
dla $x\in\R^n$, oraz
\begin{equation}\label{eq:pde4-normalize}
 \int_{\R^n}\Gamma_\kappa(t,x)\,\mathrm d\lambda_n(x)=1.
\end{equation}
Ponadto $\Gamma_\kappa(t,x)=t^{-n/2}\Gamma_\kappa(1,x/\sqrt t)$.
\end{twierdzenie}
\begin{proof}
Ustalmy $n\geq1$, $\kappa>0$ i oznaczmy
$\Gamma=\Gamma_\kappa$. Funkcja $\Gamma$ jest
dodatnia i klasy $C^\infty$ na $(0,\infty)\times\R^n$.
Różniczkując oba czynniki w jej definicji, otrzymujemy
\[
 \partial_t\Gamma(t,x)=
 \left(-\frac n{2t}+\frac{|x|^2}{4\kappa t^2}\right)
 \Gamma(t,x),\qquad
 \partial_{x_j}\Gamma(t,x)=-\frac{x_j}{2\kappa t}\Gamma(t,x).
\]
Ponowne różniczkowanie względem $x_j$ daje
\[
 \partial_{x_j}^2\Gamma(t,x)
 =\left(\frac{x_j^2}{4\kappa^2t^2}
 -\frac1{2\kappa t}\right)\Gamma(t,x).
\]
Sumując względem $j=1,\ldots,n$ i mnożąc przez
$\kappa$, dostajemy $\kappa\Delta_x\Gamma=\partial_t\Gamma$.

Dla normalizacji oznaczmy
$I=\int_\R e^{-z^2}\,\mathrm d\lambda_1(z)$.
Całka jest skończona, ponieważ dla $|z|\geq1$
mamy $e^{-z^2}\leq e^{-|z|}$.
Z twierdzenia Tonellego i zmiany zmiennych
we współrzędnych biegunowych w $\R^2$ wynika
\[
 I^2=\int_{\R^2}e^{-|z|^2}\,\mathrm d\lambda_2(z)
 =2\pi\int_0^\infty e^{-r^2}r\,\mathrm d\lambda_1(r)
 =\pi.
\]
Ponieważ $I>0$, mamy $I=\sqrt\pi$.
Ponowne zastosowanie twierdzenia Tonellego
do iloczynu $n$ funkcji jednowymiarowych daje
$\int_{\R^n}e^{-|z|^2}\,\mathrm d\lambda_n(z)=\pi^{n/2}$.
Podstawienie $x=\sqrt{4\kappa t}\,z$, którego
jakobian ma moduł $(4\kappa t)^{n/2}$, daje
\[
 \int_{\R^n}\Gamma_\kappa(t,x)\,\mathrm d\lambda_n(x)
 =(4\pi\kappa t)^{-n/2}(4\kappa t)^{n/2}\pi^{n/2}
 =1.
\]
Wreszcie wstawienie $x/\sqrt t$ do wzoru na
$\Gamma_\kappa(1,\cdot)$ daje dokładnie
$t^{n/2}\Gamma_\kappa(t,x)$.
\end{proof}

\begin{uwaga}[Skalowanie paraboliczne]
Niech $u$ będzie klasy $C^{1,2}$ na
$(0,\infty)\times\R^n$ (jedna ciągła pochodna
czasowa i ciągłe pochodne przestrzenne do rzędu
drugiego) i niech $\partial_tu-\kappa\Delta_xu=0$.
Dla $\rho>0$ połóżmy $u_\rho(t,x)=u(\rho^2t,\rho x)$.
Reguła łańcuchowa daje
\[
 (\partial_t-\kappa\Delta_x)u_\rho(t,x)
 =\rho^2(\partial_tu-\kappa\Delta_xu)(\rho^2t,\rho x)=0.
\]
Czas skaluje się zatem jak kwadrat odległości,
w odróżnieniu od skali falowej $t\sim|x|$.

Skalowanie sugeruje poszukiwanie rozwiązania o stałej
masie w postaci $v(t,x)=t^{-n/2}F(x/\sqrt t)$.
Dla $z=x/\sqrt t$ rachunek daje
\[
 (\partial_t-\kappa\Delta_x)v(t,x)
 =-t^{-n/2-1}
 \left(\kappa\Delta F(z)+\tfrac12 z\cdot\nabla F(z)
 +\tfrac n2 F(z)\right).
\]
Wyrażenie w nawiasie jest dywergencją pola
$\kappa\nabla F(z)+zF(z)/2$.
Wystarczy więc znaleźć profil spełniający
$\kappa\nabla F(z)+zF(z)/2=0$.
Dla $F\in C^1(\R^n)$ jest to równoważne
$\nabla(e^{|z|^2/(4\kappa)}F(z))=0$, czyli
$F(z)=C e^{-|z|^2/(4\kappa)}$.
Warunek jednostkowej masy daje
$C=(4\pi\kappa)^{-n/2}$.
Otrzymujemy jądro Gaussa. Nie twierdzimy, że każde
rozwiązanie równania dla profilu musi spełniać ten
silniejszy warunek zerowania pola; jest on tutaj
wystarczającą metodą konstrukcji szukanego jądra.
\end{uwaga}

\begin{center}
\begin{tikzpicture}[>=Stealth,x=1.15cm,y=3.8cm]
 \draw[->] (-3.2,0)--(3.55,0) node[right] {$x$};
 \draw[->] (0,0)--(0,1.04) node[above] {$\Gamma_1(t,x)$};
 \draw[blue!70!black,thick,domain=-3:3,samples=110]
   plot (\x,{1/sqrt(4*pi*.12)*exp(-\x*\x/(4*.12))});
 \draw[red!70!black,thick,domain=-3:3,samples=110]
   plot (\x,{1/sqrt(4*pi*.45)*exp(-\x*\x/(4*.45))});
 \draw[green!45!black,thick,domain=-3:3,samples=110]
   plot (\x,{1/sqrt(4*pi*1.1)*exp(-\x*\x/(4*1.1))});
 \node[blue!70!black,anchor=west] at (1.5,.79) {$t=0{,}12$};
 \node[red!70!black,anchor=west] at (1.5,.68) {$t=0{,}45$};
 \node[green!45!black,anchor=west] at (1.5,.57) {$t=1{,}1$};
\end{tikzpicture}
\par\smallskip\small
Jądro cieplne w wymiarze jeden ($\kappa=1$): profil
staje się niższy i szerszy, lecz całkowita masa pozostaje równa $1$.
\end{center}

\section{Własność półgrupy i przybliżona jedność}
\label{sec:pde4-semigroup}

\begin{lemat}[Składanie jąder Gaussa]
\label{lem:pde4-convolution}
Dla $n\geq1$, $\kappa>0$ oraz $s,t>0$ zachodzi
$\Gamma_\kappa(t)*\Gamma_\kappa(s)=\Gamma_\kappa(t+s)$.
\end{lemat}
\begin{proof}
Ustalmy $x\in\R^n$. Całka definiująca splot
jest skończona, ponieważ jedno jądro jest ograniczone,
a drugie ma całkę $1$. Rozwinięcie kwadratów daje
\[
 \frac{|x-y|^2}{t}+\frac{|y|^2}{s}
 =\frac{t+s}{ts}\left|y-\frac{s}{t+s}x\right|^2
 +\frac{|x|^2}{t+s}.
\]
Oznaczmy $a=ts/(t+s)>0$. Po podstawieniu
$z=y-sx/(t+s)$ otrzymujemy
\[
 \begin{aligned}
 (\Gamma_\kappa(t)*\Gamma_\kappa(s))(x)
 &=\frac{e^{-|x|^2/(4\kappa(t+s))}}
 {(4\pi\kappa)^n(ts)^{n/2}}
 \int_{\R^n}e^{-|z|^2/(4\kappa a)}
 \,\mathrm d\lambda_n(z)\\
 &=\frac{(4\pi\kappa a)^{n/2}}
 {(4\pi\kappa)^n(ts)^{n/2}}
 e^{-|x|^2/(4\kappa(t+s))}\\
 &=\frac{e^{-|x|^2/(4\kappa(t+s))}}
 {(4\pi\kappa(t+s))^{n/2}}
 =\Gamma_\kappa(t+s,x).
 \end{aligned}
\]
W drugiej równości wykorzystaliśmy normalizację
jądra w chwili $a$, a w trzeciej tożsamość
$a/(ts)=1/(t+s)$.
\end{proof}

\begin{twierdzenie}[Przybliżona jedność]
\label{thm:pde4-approx-identity}
Niech $n\geq1$, $\kappa>0$. Dla $g\in\mathrm{BUC}(\R^n)$
funkcje $P_tg:=\Gamma_\kappa(t)*g$ spełniają
$\|P_tg-g\|_\infty\to0$ przy $t\downarrow0$.
Dla $1\leq p<\infty$ i $g\in L^p(\R^n)$ zachodzi
$\|P_tg-g\|_p\to0$ przy $t\downarrow0$.
Dla każdej $g\in L^p$ ($1\leq p\leq\infty$),
$\|P_tg\|_p\leq\|g\|_p$ dla $t>0$.
\end{twierdzenie}
\begin{proof}
Dla ustalonego $t>0$ jądro $\Gamma_\kappa(t,\cdot)$
należy do każdego $L^q(\R^n)$, $1\leq q\leq\infty$:
jest ograniczone, całkowalne, a dla $q<\infty$
można oszacować jego $q$-tą potęgę przez
$\|\Gamma_\kappa(t,\cdot)\|_\infty^{q-1}$
razy samo jądro. Dla $g\in L^p(\R^n)$
nierówność Höldera z wykładnikiem sprzężonym $q$
($1/p+1/q=1$) daje więc bezwzględną zbieżność
całki definiującej $P_tg(x)$ dla każdego $x$.
Wynik nie zależy od wyboru reprezentanta klasy $g$,
ponieważ translacja zbioru miary zero ma miarę zero.

Potrzebujemy najpierw oszacowania ogona jądra.
Dla każdego $\delta>0$, korzystając ze skalowania
z twierdzenia~\ref{thm:pde4-gauss-basic}, mamy
\[
 \int_{|y|\geq\delta}\Gamma_\kappa(t,y)\,
 \mathrm d\lambda_n(y)
 =\int_{|z|\geq\delta/\sqrt t}\Gamma_\kappa(1,z)\,
 \mathrm d\lambda_n(z)\longrightarrow0
 \quad(t\downarrow0).
\]
Istotnie, funkcje charakterystyczne tych zbiorów
dążą punktowo do zera, a funkcją dominującą
jest całkowalne jądro $\Gamma_\kappa(1,\cdot)$.

Dla $g\in\mathrm{BUC}(\R^n)$ normalizacja daje
\[
 |P_tg(x)-g(x)|\leq
 \int_{\R^n}\Gamma_\kappa(t,y)|g(x-y)-g(x)|
 \,\mathrm d\lambda_n(y).
\]
Ustalmy $\eps>0$ i wybierzmy $\delta>0$ tak,
aby $|g(x-y)-g(x)|<\eps$ dla $|y|<\delta$
i wszystkich $x\in\R^n$. Dzieląc całkę na
$|y|<\delta$ oraz $|y|\geq\delta$, otrzymujemy
\[
 \|P_tg-g\|_\infty\leq\eps+
 2\|g\|_\infty\int_{|y|\geq\delta}\Gamma_\kappa(t,y)
 \,\mathrm d\lambda_n(y).
\]
Granica górna przy $t\downarrow0$ nie przekracza
$\eps$, a dowolność $\eps$ daje żądaną zbieżność.

Niech teraz $1\leq p<\infty$. Przesunięcie
$\tau_yg=g(\cdot-y)$ jest izometrią $L^p(\R^n)$.
Dla $g\in C_c^\infty(\R^n)$ zbieżność
$\|\tau_yg-g\|_p\to0$ przy $y\to0$ wynika
z jednostajnej ciągłości i wspólnego ograniczonego
nośnika dla $|y|\leq1$.
Dla dowolnego $g\in L^p$ wybieramy
$\phi\in C_c^\infty(\R^n)$ z $\|g-\phi\|_p<\eps$,
korzystając z gęstości tej klasy, i otrzymujemy
$\|\tau_yg-g\|_p\leq2\eps+\|\tau_y\phi-\phi\|_p$.
Zatem przesunięcia są ciągłe w normie $L^p$.

Nierówność całkowa Minkowskiego i normalizacja
jądra dają
\[
 \begin{aligned}
 \|P_tg\|_p&\leq\int_{\R^n}\Gamma_\kappa(t,y)
                  \|\tau_yg\|_p\,\mathrm d\lambda_n(y)=\|g\|_p,\\
 \|P_tg-g\|_p&\leq\int_{\R^n}\Gamma_\kappa(t,y)
                  \|\tau_yg-g\|_p\,\mathrm d\lambda_n(y).
 \end{aligned}
\]
Prawe strony są skończone, ponieważ
$\|\tau_yg-g\|_p\leq2\|g\|_p$.
Dla $\eps>0$ z ciągłości przesunięć wybieramy
$\delta>0$ tak, aby $\|\tau_yg-g\|_p<\eps$
dla $|y|<\delta$. Dzieląc drugą całkę jak poprzednio,
dostajemy granicę górną nie większą niż $\eps$.
Dowolność $\eps$ kończy dowód.

Dla $p=\infty$ nierówność kontrakcji wynika
bezpośrednio z $|g(x-y)|\leq\|g\|_\infty$
dla prawie wszystkich $y$ oraz normalizacji jądra.
\end{proof}

Ustalmy $\kappa>0$ i oznaczmy $P_0g=g$.
Dla $s,t>0$ i $g\in L^p$, $1\leq p\leq\infty$,
podwójna całka używana w $P_tP_sg(x)$ jest
bezwzględnie zbieżna. Istotnie, twierdzenie Tonellego
dla całki z $|g|$ oraz lemat~\ref{lem:pde4-convolution}
dają
\[
 \int_{\R^n}\int_{\R^n}\Gamma_\kappa(t,y)
 \Gamma_\kappa(s,z)|g(x-y-z)|\,
 \mathrm d\lambda_n(z)\,\mathrm d\lambda_n(y)
 =P_{t+s}|g|(x)<\infty.
\]
Można zatem zastosować twierdzenie Fubiniego
i otrzymać $P_tP_sg=P_{t+s}g$. Ta sama równość
zachodzi dla $s=0$ lub $t=0$ z definicji $P_0$.

Jeśli $g\in\mathrm{BUC}(\R^n)$, to
\[
 |P_tg(x+h)-P_tg(x)|
 \leq\sup_{z\in\R^n}|g(z+h)-g(z)|,
\]
co pokazuje, że $P_t$ zachowuje $\mathrm{BUC}$
i jest w niej kontrakcją.
W przestrzeni $X=L^p(\R^n)$ dla $1\leq p<\infty$
albo $X=\mathrm{BUC}(\R^n)$ dla $g\in X$ mamy
\[
 \|P_tg-P_sg\|_X\leq\|P_{t-s}g-g\|_X
 \qquad(0\leq s\leq t).
\]
Wynika to z własności półgrupy i kontrakcji
$P_s$. Twierdzenie~\ref{thm:pde4-approx-identity}
daje zatem ciągłość odwzorowania $t\mapsto P_tg$
w normie $X$ na $[0,\infty)$, również z lewej
strony w każdym punkcie dodatnim.

Na $L^\infty(\R^n)$ silna ciągłość w chwili zero
nie musi zachodzić. Istotnie, niech
$g(x)=\chi_{\{x_1>0\}}(x)$.
Z symetrii jądra względem zmiany znaku pierwszej
współrzędnej mamy $P_tg(0)=1/2$ dla każdego $t>0$.
Ponadto
\[
 P_tg(x)=\int_{\{y_1<x_1\}}\Gamma_\kappa(t,y)\,
 \mathrm d\lambda_n(y),
\]
więc funkcja po lewej jest ciągła w $x$ dzięki
twierdzeniu o zbieżności zdominowanej (hiperpłaszczyzna
$y_1=x_1$ ma miarę Lebesgue'a zero).
Dla każdego $\eta>0$ na pewnej półkuli
$B_\delta(0)\cap\{x_1>0\}$ dodatniej miary
zachodzi $|P_tg(x)-1/2|<\eta$ i $g(x)=1$.
Stąd $\|P_tg-g\|_\infty\geq1/2-\eta$,
a dowolność $\eta$ daje
$\|P_tg-g\|_\infty\geq1/2$ dla każdego $t>0$.
Przestrzeń $\mathrm{BUC}$ jest więc naturalną
przestrzenią danych ograniczonych dla silnie
ciągłej półgrupy cieplnej.

\section{Zagadnienie Cauchy'ego na całej przestrzeni}
\label{sec:pde4-cauchy}

\begin{twierdzenie}[Jawne rozwiązanie z danych początkowych]
\label{thm:pde4-cauchy}
Niech $n\geq1$, $\kappa>0$ i $g\in\mathrm{BUC}(\R^n)$.
Definiujemy dla $t>0$,
\begin{equation}\label{eq:pde4-cauchy}
 u(t,x)=\int_{\R^n}\Gamma_\kappa(t,x-y)g(y)
 \,\mathrm d\lambda_n(y),\qquad u(0,x)=g(x).
\end{equation}
Wówczas $u\in C([0,\infty);\mathrm{BUC}(\R^n))$,
$u\in C^\infty((0,\infty)\times\R^n)$,
$u_t-\kappa\Delta_xu=0$ dla $t>0$ oraz
$\sup_{t\geq0}\|u(t,\cdot)\|_\infty\leq\|g\|_\infty$.
Ponadto $g\geq0$ implikuje $u\geq0$.
Jest to jedyne rozwiązanie z klasy
$C^{1,2}((0,T]\times\R^n)\cap C([0,T]\times\R^n)$,
ograniczone na $[0,T]\times\R^n$, o tych danych,
dla każdego ustalonego $T>0$.
\end{twierdzenie}
\begin{proof}
Ustalmy $\kappa>0$ i zapiszmy $\Gamma=\Gamma_\kappa$.
Dla każdego $t>0$ całka w~\eqref{eq:pde4-cauchy}
jest bezwzględnie zbieżna, gdyż $g$ jest ograniczona,
a jądro ma całkę $1$. Wykażmy gładkość rozwiązania
we wszystkich zmiennych dla dodatnich czasów.

Dla każdej liczby całkowitej $m\geq0$ oraz
wielowskaźnika $\alpha\in\mathbb N_0^n$ pochodna
$\partial_t^mD_z^\alpha\Gamma(t,z)$ jest postaci
$A_{m,\alpha}(t,z)\Gamma(t,z)$, gdzie $A_{m,\alpha}$
jest wielomianem względem $z$ o współczynnikach
gładkich dla $t>0$. Istotnie, zaczynamy od $A_{0,0}=1$.
Różniczkowanie wyrażenia $A\Gamma$ daje odpowiednio
\[
 \begin{aligned}
 \partial_{z_j}(A\Gamma)
 &=\left(\partial_{z_j}A-\frac{z_j}{2\kappa t}A\right)\Gamma,\\
 \partial_t(A\Gamma)
 &=\left(\partial_tA+
 \left(-\frac n{2t}+\frac{|z|^2}{4\kappa t^2}\right)A\right)\Gamma.
 \end{aligned}
\]
Obie operacje zachowują opisaną postać, co dowodzi
jej dla wszystkich pochodnych przez indukcję.

Ustalmy $0<a<b<\infty$ i $L>0$ oraz rozważmy
$(t,x)\in[a,b]\times\overline{B_L(0)}$.
Dla ustalonej pochodnej istnieją stała $C>0$
i liczba całkowita $d\geq0$, niezależne od $t,x,y$,
takie, że $|A_{m,\alpha}(t,x-y)|\leq C(1+|x-y|)^d$.
Wynika to z ograniczoności skończonej liczby
współczynników wielomianu na $[a,b]$.
Ponadto z nierówności
$|y|^2\leq2|x-y|^2+2|x|^2$ dostajemy
$|x-y|^2\geq |y|^2/2-L^2$. W konsekwencji
\[
 e^{-|x-y|^2/(4\kappa t)}
 \leq e^{-|x-y|^2/(4\kappa b)}
 \leq e^{L^2/(4\kappa b)}e^{-|y|^2/(8\kappa b)}.
\]
Ponieważ $(4\pi\kappa t)^{-n/2}\leq(4\pi\kappa a)^{-n/2}$
i $|x-y|\leq L+|y|$, istnieje stała $C'>0$ taka,
że na rozważanym zbiorze
\[
 |\partial_t^mD_x^\alpha\Gamma(t,x-y)g(y)|
 \leq C'\|g\|_\infty(1+|y|)^d
       e^{-|y|^2/(8\kappa b)}.
\]
Prawa strona jest całkowalna w $\R^n$, ponieważ
każdy wielomian wzrostu promieniowego jest
całkowalny po pomnożeniu przez taką funkcję Gaussa.

Na wnętrzu wskazanego cylindra ilorazy różnicowe
są zdominowane przez oszacowanie odpowiedniej
pochodnej następnego rzędu, na mocy twierdzenia
o wartości średniej. Twierdzenie o zbieżności
zdominowanej uzasadnia zarówno różniczkowanie
pod całką, jak i ciągłość otrzymanych pochodnych.
Dowolność $a,b,L$ daje $u\in C^\infty((0,\infty)\times\R^n)$
i dla $t>0$ wzór
\[
 \partial_t^mD_x^\alpha u(t,x)
 =\int_{\R^n}\partial_t^mD_x^\alpha\Gamma(t,x-y)
 g(y)\,\mathrm d\lambda_n(y).
\]
Stąd, stosując twierdzenie~\ref{thm:pde4-gauss-basic},
otrzymujemy
$(\partial_t-\kappa\Delta_x)u(t,x)=0$.

Z poprzedniej sekcji wiemy, że $P_t$ zachowuje
przestrzeń $\mathrm{BUC}$ oraz że $t\mapsto P_tg$
jest ciągłe w jej normie na $[0,\infty)$,
gdzie $P_0g=g$. W szczególności $u(t,\cdot)\to g$
jednostajnie na $\R^n$ przy $t\downarrow0$.
Uzasadnijmy również ciągłość łączną: dla $s\geq0$,
$x_0\in\R^n$ i $(t,x)\to(s,x_0)$, $t\geq0$, mamy
\[
 |u(t,x)-u(s,x_0)|\leq
 \|P_tg-P_sg\|_\infty+
 |P_sg(x)-P_sg(x_0)|\longrightarrow0.
\]
Pierwszy składnik znika dzięki ciągłości normowej,
a drugi dzięki ciągłości funkcji $P_sg$.

Kontrakcyjność daje $\|u(t,\cdot)\|_\infty\leq\|g\|_\infty$
również dla $t=0$. Jeżeli $g\geq0$, to $u\geq0$
dzięki dodatniości jądra. Na każdym przedziale
$[0,T]$ skonstruowana funkcja należy więc do klasy
wymaganej w tezie o jednoznaczności.
Twierdzenie~\ref{thm:pde4-global-uniqueness},
udowodnione poniżej przy użyciu zasady maksimum,
zapewnia jednoznaczność w tej klasie. Jego dowód
nie korzysta z obecnej reprezentacji całkowej.
\end{proof}

\begin{wniosek}[Natychmiastowe rozchodzenie się zaburzenia]
\label{cor:pde4-positive}
Niech $n\geq1$, $\kappa>0$, a
$g\in C_c(\R^n)$ będzie nieujemna i nierówna identycznie
zeru. Wówczas rozwiązanie~\eqref{eq:pde4-cauchy}
ma własność $u(t,x)>0$ dla każdego $t>0$, $x\in\R^n$.
\end{wniosek}
\begin{proof}
Funkcja ciągła o zwartym nośniku jest ograniczona
i jednostajnie ciągła na $\R^n$, więc spełnia
założenia twierdzenia~\ref{thm:pde4-cauchy}.
Ponieważ $g$ jest nieujemna i nie jest identycznie
zerowa, istnieje $y_0\in\R^n$ takie, że $g(y_0)>0$.
Z ciągłości wybieramy $r>0$ i $c=g(y_0)/2>0$
tak, aby $g\geq c$ na $\overline{B_r(y_0)}$.
Dla ustalonego $t>0$ i $x\in\R^n$ funkcja
$y\mapsto\Gamma_\kappa(t,x-y)$ jest ciągła
i ściśle dodatnia, a więc osiąga dodatnie minimum
$m_{t,x}$ na tej zwartej kuli.
Nieujemność $g$ na całej przestrzeni daje
\[
 u(t,x)\geq c\int_{B_r(y_0)}
 \Gamma_\kappa(t,x-y)\,\mathrm d\lambda_n(y)
 \geq c\,m_{t,x}\lambda_n(B_r(y_0))>0.
\]
W szczególności zwarty nośnik nieujemnych,
niezerowych danych początkowych nie pociąga
za sobą ograniczonego nośnika rozwiązania
w żadnej dodatniej chwili. Nie twierdzimy tu,
że dane zmieniające znak dają rozwiązanie dodatnie.
\end{proof}

\begin{stwierdzenie}[Zachowanie masy i tłumienie norm]
\label{prop:pde4-mass}
Niech $n\geq1$, $\kappa>0$, $t>0$.
Dla $g\in L^1(\R^n)$ zachodzi
\begin{equation}\label{eq:pde4-mass}
 \int_{\R^n}(P_tg)(x)\,\mathrm d\lambda_n(x)
 =\int_{\R^n}g(x)\,\mathrm d\lambda_n(x).
\end{equation}
Dla $g\in L^1(\R^n)$ zachodzi również
\begin{equation}\label{eq:pde4-L1Linf}
 \|P_tg\|_\infty\leq (4\pi\kappa t)^{-n/2}\|g\|_1,
 \qquad
 \|P_tg\|_2\leq (8\pi\kappa t)^{-n/4}\|g\|_1.
\end{equation}
\end{stwierdzenie}
\begin{proof}
Niech $g\in L^1(\R^n)$ i ustalmy $t>0$.
Twierdzenie Tonellego dla nieujemnej funkcji
podcałkowej oraz normalizacja jądra dają
\[
 \begin{aligned}
 &\int_{\R^n}\int_{\R^n}\Gamma_\kappa(t,x-y)|g(y)|
       \,\mathrm d\lambda_n(y)\,\mathrm d\lambda_n(x)\\
 &\quad=\int_{\R^n}|g(y)|
   \left(\int_{\R^n}\Gamma_\kappa(t,x-y)\,
                 \mathrm d\lambda_n(x)\right)\mathrm d\lambda_n(y)
 =\|g\|_1.
 \end{aligned}
\]
Wewnętrzna całka ma wartość $1$ dzięki niezmienniczości
miary Lebesgue'a względem przesunięć.
W szczególności podwójna całka z $g$ zamiast
$|g|$ jest bezwzględnie zbieżna. Stosując
twierdzenie Fubiniego otrzymujemy
\[
 \int_{\R^n}P_tg(x)\,\mathrm d\lambda_n(x)
 =\int_{\R^n}g(y)\,\mathrm d\lambda_n(y),
\]
czyli~\eqref{eq:pde4-mass}. Jest to zachowanie
całki z funkcji, również dla danych zmieniających
znak. Jeśli $g\geq0$ prawie wszędzie, dodatniość
jądra daje $P_tg\geq0$ i stąd
$\|P_tg\|_1=\|g\|_1$. Dla dowolnych danych
$|P_tg|\leq P_t|g|$, co daje tylko kontrakcję
w normie $L^1$.

Maksimum jądra jest osiągane w zerze:
$\Gamma_\kappa(t,0)=(4\pi\kappa t)^{-n/2}$.
Zatem dla każdego $x\in\R^n$ zachodzi
\[
 |P_tg(x)|\leq(4\pi\kappa t)^{-n/2}\|g\|_1.
\]
Biorąc istotne supremum po $x$, otrzymujemy
pierwszą nierówność w~\eqref{eq:pde4-L1Linf}.

Dla drugiej obliczamy normę jądra wprost.
Korzystając z całki Gaussa, mamy
\[
 \begin{aligned}
 \|\Gamma_\kappa(t,\cdot)\|_2^2
 &=(4\pi\kappa t)^{-n}
 \int_{\R^n}e^{-|x|^2/(2\kappa t)}\,\mathrm d\lambda_n(x)\\
 &=(4\pi\kappa t)^{-n}(2\pi\kappa t)^{n/2}
 =(8\pi\kappa t)^{-n/2}.
 \end{aligned}
\]
Z niezmienniczości normy względem przesunięć
i całkowej nierówności Minkowskiego wynika
\[
 \begin{aligned}
 \|P_tg\|_2
 &\leq\int_{\R^n}|g(y)|
       \|\Gamma_\kappa(t,\cdot-y)\|_2\,\mathrm d\lambda_n(y)\\
 &= (8\pi\kappa t)^{-n/4}\|g\|_1.
 \end{aligned}
\]
Prawa strona jest skończona dla $g\in L^1$,
co zarazem uzasadnia, że $P_tg\in L^2$.
\end{proof}

\section{Paraboliczna zasada maksimum i jednoznaczność}
\label{sec:pde4-maximum}

Niech $\Omega\subset\R^n$ będzie obszarem ograniczonym,
a $T>0$. Wprowadzamy cylinder $Q_T=(0,T)\times\Omega$,
jego domknięcie $\overline{Q_T}=[0,T]\times\overline\Omega$
oraz \emph{brzeg paraboliczny}
\[
 \partial_pQ_T=(\{0\}\times\overline\Omega)
 \cup([0,T]\times\partial\Omega).
\]
Brzeg paraboliczny obejmuje dolną podstawę i ścianę
boczną cylindra, ale nie obejmuje punktów
$\{T\}\times\Omega$. To właśnie na brzegu parabolicznym
kontrolujemy dane rozwiązania w zagadnieniu postępowym;
na górnej podstawie nie zadaje się niezależnych danych.

Przez $C^{1,2}((0,T]\times\Omega)$ rozumiemy tutaj
funkcje o ciągłej pochodnej czasowej i ciągłych
pochodnych przestrzennych do rzędu drugiego
na wskazanym zbiorze; w chwili $T$ pochodna czasowa
jest rozumiana jako pochodna lewostronna.
Nie zakładamy istnienia pochodnych na brzegu
przestrzennym ani w chwili początkowej.

\begin{twierdzenie}[Słaba paraboliczna zasada maksimum]
\label{thm:pde4-max}
Niech $n\geq1$, $\Omega\subset\R^n$ będzie obszarem
ograniczonym, $T>0$, $\kappa>0$, a
$u\in C([0,T]\times\overline\Omega)\cap
C^{1,2}((0,T]\times\Omega)$ będzie funkcją rzeczywistą
spełniającą
$\partial_tu-\kappa\Delta_xu\leq0$
na $(0,T]\times\Omega$.
Wówczas
\begin{equation}\label{eq:pde4-max}
 \max_{[0,T]\times\overline\Omega}u
 \leq \max_{\partial_pQ_T}u.
\end{equation}
W konsekwencji, jeśli funkcje $u_1,u_2$ z tej samej
klasy spełniają
$(\partial_t-\kappa\Delta_x)u_1
\leq(\partial_t-\kappa\Delta_x)u_2$
na $(0,T]\times\Omega$ oraz $u_1\leq u_2$
na $\partial_pQ_T$, to $u_1\leq u_2$
na całym $\overline{Q_T}$.
W szczególności rozwiązania tego samego równania
z tymi samymi danymi na brzegu parabolicznym
są jednoznaczne w podanej klasie.
\end{twierdzenie}
\begin{proof}
Ustalmy $\eps>0$ i połóżmy $v(t,x)=u(t,x)-\eps t$.
Funkcja $v$ ma tę samą regularność co $u$ oraz spełnia
\[
 \partial_tv-\kappa\Delta_xv
 =\partial_tu-\kappa\Delta_xu-\eps\leq-\eps<0.
\]
Ponieważ $\Omega$ jest ograniczony,
$\overline{Q_T}$ jest zwarty. Funkcja $v$ osiąga
więc na nim maksimum. Przypuśćmy, że pewien
punkt maksimum $(t_0,x_0)$ nie należy do
$\partial_pQ_T$. Wówczas $0<t_0\leq T$ i $x_0\in\Omega$.
Funkcja $x\mapsto v(t_0,x)$ osiąga maksimum
lokalne w punkcie wewnętrznym $x_0$. Stąd
$\nabla_xv(t_0,x_0)=0$, jej hesjan przestrzenny
jest niedodatnio określony i w szczególności
$\Delta_xv(t_0,x_0)\leq0$.

Jeśli $t_0<T$, funkcja $t\mapsto v(t,x_0)$
ma maksimum w punkcie wewnętrznym przedziału
$(0,T)$, a więc $\partial_tv(t_0,x_0)=0$.
Jeśli natomiast $t_0=T$, dla każdego dostatecznie
małego $h>0$ mamy
\[
 \frac{v(T,x_0)-v(T-h,x_0)}{h}\geq0.
\]
Przechodząc z $h$ do zera otrzymujemy
$\partial_tv(T,x_0)\geq0$ w sensie pochodnej
lewostronnej. W obu przypadkach
$(\partial_t-\kappa\Delta_x)v(t_0,x_0)\geq0$,
co przeczy ścisłej nierówności dla $v$.

Każdy punkt maksimum $v$ należy zatem do brzegu
parabolicznego. Na tym brzegu $v\leq u$, gdyż
$t\geq0$, a na całym domkniętym cylindrze
$u=v+\eps t\leq v+\eps T$. Otrzymujemy
\[
 \max_{\overline{Q_T}}u
 \leq\max_{\overline{Q_T}}v+\eps T
 =\max_{\partial_pQ_T}v+\eps T
 \leq\max_{\partial_pQ_T}u+\eps T.
\]
Dowolność $\eps>0$ daje~\eqref{eq:pde4-max}.
Nie użyliśmy żadnego założenia o gładkości
brzegu przestrzennego $\Omega$.

Dla porównania funkcji $u_1,u_2$ stosujemy
wykazaną nierówność do $w=u_1-u_2$.
Z założeń wynika
$(\partial_t-\kappa\Delta_x)w\leq0$
oraz $w\leq0$ na brzegu parabolicznym, więc $w\leq0$
na $\overline{Q_T}$. Jeżeli obie funkcje spełniają
to samo równanie i mają te same dane, stosujemy
porównanie także w odwrotnym kierunku.
\end{proof}

\begin{twierdzenie}[Jednoznaczność w klasie rozwiązań ograniczonych]
\label{thm:pde4-global-uniqueness}
Niech $n\geq1$, $\kappa>0$, $T>0$. Jeśli
$u,v\in C([0,T]\times\R^n)\cap
C^{1,2}((0,T]\times\R^n)$ są rzeczywistymi
funkcjami ograniczonymi na $[0,T]\times\R^n$,
spełniają
\[
 \partial_tu-\kappa\Delta_xu=0,\qquad
 \partial_tv-\kappa\Delta_xv=0
 \quad\text{na }(0,T]\times\R^n,
\]
to zachodzi oszacowanie stabilności względem
danych początkowych
\[
 \sup_{(t,x)\in[0,T]\times\R^n}|u(t,x)-v(t,x)|
 \leq\sup_{x\in\R^n}|u(0,x)-v(0,x)|.
\]
W szczególności, jeśli $u(0,x)=v(0,x)$ dla wszystkich
$x\in\R^n$, to $u=v$ na $[0,T]\times\R^n$.
\end{twierdzenie}
\begin{proof}
Niech $w=u-v$ i oznaczmy
$M=\sup_{[0,T]\times\R^n}|w|<\infty$ oraz
$A=\sup_{\R^n}|w(0,\cdot)|\leq M$.
Funkcja $w$ spełnia jednorodne równanie ciepła.
Nie możemy jednak bezpośrednio zastosować
twierdzenia~\ref{thm:pde4-max} do całej przestrzeni,
gdyż cylinder jest nieograniczony.

Dla $\eps>0$ wprowadzamy funkcję
$b_\eps(t,x)=\eps(|x|^2+2n\kappa t)$.
Ponieważ $\Delta_x|x|^2=2n$, zachodzi
\[
 \partial_tb_\eps=2n\kappa\eps
 =\kappa\Delta_xb_\eps.
\]
Funkcja $b_\eps$ jest więc dokładnym rozwiązaniem
równania ciepła. Jest nieujemna dla $t\geq0$
i rośnie kwadratowo wraz z $|x|$.
Ustalmy dowolny punkt $(t_*,x_*)\in[0,T]\times\R^n$.
Dla danego $\eps>0$ wybierzmy $R>0$ tak, aby
$R>|x_*|$ i $\eps R^2>M$.

Na cylindrze $[0,T]\times\overline{B_R(0)}$
rozważamy $w_\eps=w-A-b_\eps$.
Funkcja ta jest ciągła na domknięciu cylindra,
ma wymaganą regularność wewnątrz i na górnej
podstawie oraz spełnia
$(\partial_t-\kappa\Delta_x)w_\eps=0$.
Na dolnej podstawie mamy
\[
 w_\eps(0,x)=w(0,x)-A-\eps|x|^2\leq0
 \qquad(x\in\overline{B_R(0)}).
\]
Na ścianie bocznej dla $0\leq t\leq T$ i $|x|=R$
otrzymujemy
$w_\eps(t,x)\leq M-A-\eps R^2<0$.
Twierdzenie~\ref{thm:pde4-max} daje więc
$w_\eps\leq0$ na całym domkniętym cylindrze.
W szczególności
\[
 w(t_*,x_*)\leq A+\eps(|x_*|^2+2n\kappa t_*).
\]
Chociaż promień $R$ zależy od $\eps$, powyższe
oszacowanie w ustalonym punkcie zachodzi dla
każdego $\eps>0$. Przechodząc z $\eps$ do zera,
otrzymujemy $w(t_*,x_*)\leq A$.

Stosujemy ten sam argument do $-w$, dla którego
stałe $M$ i $A$ są takie same. Otrzymujemy
$-w(t_*,x_*)\leq A$. Punkt był dowolny, zatem
$|w|\leq A$ na $[0,T]\times\R^n$, co dowodzi
oszacowania. Dla zgodnych danych początkowych
$A=0$, więc $w=0$.
Dowód nie korzysta z reprezentacji rozwiązania
za pomocą splotu z jądrem Gaussa.
\end{proof}

\begin{uwaga}[Dlaczego trzeba podać klasę rozwiązań?]
Ograniczoność na całym $[0,T]\times\R^n$ jest
w powyższym dowodzie istotna: pozwala dla danego
$\eps$ wybrać promień kuli niezależny od czasu,
na której ścianie bocznej dominuje bariera.
Nie należy utożsamiać tego założenia z samą
ograniczonością na zbiorach zwartych, która
wynika już z ciągłości rozwiązania.
Pod koniec rozdziału przedstawimy niezerowe
rozwiązanie jednorodnego równania ciepła z zerowymi
danymi początkowymi, niespełniające globalnego
warunku ograniczoności.
\end{uwaga}

\section{Wygładzanie i analityczność przestrzenna}
\label{sec:pde4-smoothing}

\begin{twierdzenie}[Oszacowania pochodnych]
\label{thm:pde4-smoothing}
Niech $n\geq1$, $\kappa>0$, $g\in L^\infty(\R^n)$,
a $u(t,x)=P_tg(x)$ dla $t>0$.
Dla każdego wielowskaźnika $\alpha\in\mathbb N_0^n$
istnieje stała $C_{\alpha,n}>0$, niezależna od
$g,t,\kappa$, taka że dla każdego $t>0$
\begin{equation}\label{eq:pde4-derivative}
 \|\partial_x^\alpha u(t,\cdot)\|_\infty
 \leq C_{\alpha,n}(\kappa t)^{-|\alpha|/2}\|g\|_\infty.
\end{equation}
Co więcej, dla każdego ustalonego $t>0$ funkcja
$x\mapsto u(t,x)$ przedłuża się do funkcji holomorficznej
$U_t$ na $\C^n$, spełniającej
\[
 |U_t(x+i\eta)|\leq
 e^{|\eta|^2/(4\kappa t)}\|g\|_\infty
 \qquad(x,\eta\in\R^n).
\]
\end{twierdzenie}
\begin{proof}
Oznaczmy $a=\kappa t$ i $K(z)=(4\pi)^{-n/2}e^{-|z|^2/4}$
dla $z\in\R^n$. Wówczas $\Gamma_\kappa(t,x)=a^{-n/2}K(x/\sqrt a)$.
Dla wielowskaźnika $\alpha=(\alpha_1,\ldots,\alpha_n)$
o długości $|\alpha|=\sum_j\alpha_j$ reguła łańcuchowa daje
\[
 \partial_x^\alpha\Gamma_\kappa(t,x)
 =a^{-(n+|\alpha|)/2}(\partial^\alpha K)(x/\sqrt a).
\]
Ponieważ $\partial_{z_j}K=-z_jK/2$, indukcyjnie otrzymujemy
$\partial^\alpha K=Q_\alpha K$, gdzie $Q_0=1$ oraz
$Q_{\alpha+e_j}=\partial_{z_j}Q_\alpha-z_jQ_\alpha/2$;
$e_j$ jest $j$-tym wektorem standardowej bazy.
Wielomian $Q_\alpha$ ma niezerowy składnik jednorodny
najwyższego stopnia $(-1/2)^{|\alpha|}z^\alpha$.
Zatem $C_{\alpha,n}=\|\partial^\alpha K\|_1$ jest
skończone i dodatnie. Podstawienie $x=\sqrt a\,z$
daje $\|\partial_x^\alpha\Gamma_\kappa(t,\cdot)\|_1
=a^{-|\alpha|/2}C_{\alpha,n}$.
Argument dominacji z dowodu twierdzenia
\ref{thm:pde4-cauchy} wymaga tylko istotnej
ograniczoności $g$, nie jej ciągłości: wybieramy
ograniczonego reprezentanta klasy. Wolno więc
różniczkować splot dowolną liczbę razy dla $t>0$.
Z niezmienniczości miary Lebesgue'a względem przesunięć
otrzymujemy
\[
 |\partial_x^\alpha u(t,x)|
 \leq\|g\|_\infty\int_{\R^n}
 |\partial_x^\alpha\Gamma_\kappa(t,x-y)|\,\mathrm d\lambda_n(y)
 =a^{-|\alpha|/2}C_{\alpha,n}\|g\|_\infty.
\]
Biorąc supremum otrzymujemy~\eqref{eq:pde4-derivative}.

Dla ustalonego $t>0$ definiujemy
\[
 U_t(z)=\int_{\R^n}E(z,y)g(y)\,\mathrm d\lambda_n(y),
 \qquad
 E(z,y)=\frac1{(4\pi a)^{n/2}}
 \exp\!\left(-\frac{\sum_{j=1}^n(z_j-y_j)^2}{4a}\right).
\]
W wykładniku występuje suma kwadratów zespolonych,
a nie kwadrat normy hermitowskiej. Dla $z=x+i\eta$
mamy $\operatorname{Re}\sum_j(z_j-y_j)^2=|x-y|^2-|\eta|^2$,
więc $|E(x+i\eta,y)|=e^{|\eta|^2/(4a)}\Gamma_\kappa(t,x-y)$.
Normalizacja jądra daje bezwzględną zbieżność
całki dla każdego $z$ oraz oszacowanie podane w tezie.

Dla każdego $y$ funkcja $E(\cdot,y)$ jest holomorficzna
na $\C^n$. Jeżeli $|\operatorname{Re}z|,|\operatorname{Im}z|\leq L$,
to z $|x-y|^2\geq |y|^2/2-L^2$ wynika
\[
 |E(z,y)|\leq(4\pi a)^{-n/2}e^{L^2/(2a)}e^{-|y|^2/(8a)}.
\]
Ponadto
\[
 \begin{aligned}
 \partial_{z_j}E&=-\frac{z_j-y_j}{2a}E,\\
 \partial_{z_k}\partial_{z_j}E
 &=\left(\frac{(z_j-y_j)(z_k-y_k)}{4a^2}
 -\frac{\delta_{jk}}{2a}\right)E,
 \end{aligned}
\]
gdzie $\delta_{jk}=1$ dla $j=k$ i $0$ w przeciwnym
przypadku. Funkcję $E$ i obie grupy pochodnych,
pomnożone przez $g(y)$, dominuje na wskazanym
zbiorze jedna całkowalna funkcja postaci
$C\|g\|_\infty(1+|y|)^2e^{-|y|^2/(8a)}$.

Ustalmy $z_0\in\C^n$ i wybierzmy $L$ tak duże,
aby odcinki $z_0+sh$, $0\leq s\leq1$, $|h|\leq1$,
mieściły się w rozważanym zbiorze. Wzór Taylora
z resztą całkową dla funkcji $s\mapsto E(z_0+sh,y)$ daje
\[
 \begin{aligned}
 E(z_0+h,y)-E(z_0,y)
 &=\sum_j\partial_{z_j}E(z_0,y)h_j\\
 &\quad+\int_0^1(1-s)\sum_{j,k}
 \partial_{z_j}\partial_{z_k}E(z_0+sh,y)h_jh_k\,
 \mathrm d\lambda_1(s).
 \end{aligned}
\]
Ponieważ $\sum_{j,k}|h_jh_k|\leq n|h|^2$,
moduł reszty pomnożonej przez $g$ jest ograniczony
przez całkowalną majorantę razy $|h|^2$. Po całkowaniu
otrzymujemy
\[
 U_t(z_0+h)=U_t(z_0)+\sum_jh_j
 \int_{\R^n}\partial_{z_j}E(z_0,y)g(y)\,\mathrm d\lambda_n(y)
 +O(|h|^2).
\]
Wyraz liniowy jest zespolono-liniowy względem $h$.
Funkcja $U_t$ jest różniczkowalna w sensie Frécheta
nad $\C$ w każdym punkcie $\C^n$, a zatem holomorficzna.
Jej ograniczenie do $\R^n$ jest równe $P_tg$,
w szczególności jest rzeczywiście analityczne.
\end{proof}

\begin{przyklad}[Wygasanie wysokich częstości]
Niech $\xi\in\R^n$ i $g(x)=\cos(\xi\cdot x)$.
Funkcja $g$ należy do $\mathrm{BUC}(\R^n)$,
gdyż jest ograniczona i ma ograniczony gradient.
Połóżmy $u(t,x)=e^{-\kappa|\xi|^2t}\cos(\xi\cdot x)$.
Mamy $\partial_tu=-\kappa|\xi|^2u$ oraz
$\partial_{x_j}^2u=-\xi_j^2u$, więc
$\partial_tu-\kappa\Delta_xu=0$. Funkcja jest gładka,
spełnia dane początkowe i jest ograniczona przez $1$
na całym $[0,\infty)\times\R^n$. Z twierdzeń
\ref{thm:pde4-cauchy} i~\ref{thm:pde4-global-uniqueness}
wynika $u(t,\cdot)=P_tg$. Dla ustalonego $t>0$
mnożnik $e^{-\kappa t|\xi|^2}$ maleje wykładniczo
względem kwadratu normy częstości przestrzennej.
Dla $\xi=0$ rozwiązanie jest stałe i nie zanika.
\end{przyklad}

\begin{lemat}[Generator półgrupy cieplnej na gładkich danych]
\label{lem:pde-v31-heat-generator}
Niech $n\geq1$, $\kappa>0$ oraz
$h\in C^2(\R^n)$, przy czym $h$ i wszystkie pochodne
przestrzenne rzędu co najwyżej dwa są ograniczone
i jednostajnie ciągłe na $\R^n$.
Wtedy dla $t>0$ zachodzi
\begin{equation}\label{eq:pde-v31-heat-generator}
 P_th-h=\kappa\int_0^t P_s(\Delta h)\,\mathrm d\lambda_1(s),
 \qquad
 \lim_{t\downarrow0}\left\|\frac{P_th-h}{t}
                 -\kappa\Delta h\right\|_\infty=0.
\end{equation}
Całkę rozumiemy punktowo lub, równoważnie, jako
całkę funkcji ciągłej o wartościach w $\mathrm{BUC}(\R^n)$.
\end{lemat}
\begin{proof}
Dla $s>0$ mamy $P_sh(x)=\int_{\R^n}\Gamma_\kappa(s,z)
h(x-z)\,\mathrm d\lambda_n(z)$. Różniczkując pod
całką względem $x$ do rzędu drugiego, korzystamy
z ograniczoności pochodnych $h$ i całkowalności
jądra: ilorazy różnicowe są zdominowane przez
jądro pomnożone przez supremum odpowiedniej pochodnej $h$.
Otrzymujemy $\Delta_xP_sh=P_s(\Delta h)$.
Z twierdzenia~\ref{thm:pde4-cauchy} wynika więc
$\partial_sP_sh=\kappa\Delta_xP_sh=\kappa P_s(\Delta h)$.

Alternatywnie można przenieść laplasjan z jądra
na dane. Ustalmy $s>0$, $x\in\R^n$ i oznaczmy
$K_x(y)=\Gamma_\kappa(s,x-y)$. Mamy
$\Delta_x K_x=\Delta_y K_x$. Całkowanie przez części
na $B_R(0)$ daje
\[
 \int_{B_R}(h\Delta_y K_x-K_x\Delta h)\,\mathrm d\lambda_n
 =\int_{\partial B_R}(h\partial_\nu K_x-K_x\partial_\nu h)
 \,\mathrm d\sigma.
\]
Dla $R>|x|$ moduł całki brzegowej jest nie większy niż
\[
 \begin{aligned}
 &\sigma(S^{n-1})R^{n-1}(4\pi\kappa s)^{-n/2}
 e^{-(R-|x|)^2/(4\kappa s)}\\
 &\qquad{}\cdot\left(\|h\|_\infty\frac{R+|x|}{2\kappa s}
 +\|\nabla h\|_\infty\right).
 \end{aligned}
\]
W wymiarze jeden przyjmujemy $\sigma(S^0)=2$,
a całka brzegowa jest sumą po końcach przedziału.
Oszacowanie dąży do zera przy $R\to\infty$.
Całki objętościowe są bezwzględnie zbieżne,
gdyż $h,\Delta h$ są ograniczone, a jądro
i jego pochodne są całkowalne. Przejście do granicy
daje tę samą tożsamość $\Delta_xP_sh=P_s(\Delta h)$.

Dla $0<\delta<t$ całkujemy punktowo równanie
względem czasu, otrzymując
\[
 P_th(x)-P_\delta h(x)
 =\kappa\int_\delta^t P_s(\Delta h)(x)\,\mathrm d\lambda_1(s).
\]
Ponieważ $\Delta h\in\mathrm{BUC}$, odwzorowanie
$s\mapsto P_s(\Delta h)$ jest ciągłe w normie
na $[0,t]$. Jego całka w tej przestrzeni istnieje
jako granica sum Riemanna; ciągłość funkcjonałów
obliczania wartości w punkcie daje zgodność z całką
punktową. Mamy
$\|\int_0^\delta P_s(\Delta h)\,\mathrm d\lambda_1(s)\|_\infty
\leq\delta\|\Delta h\|_\infty$ oraz
$P_\delta h\to h$ w normie supremum.
Przejście z $\delta$ do zera daje pierwszą równość.

Dzieląc ją przez $t$ i odejmując $\kappa\Delta h$,
otrzymujemy
\[
 \left\|\frac{P_th-h}{t}-\kappa\Delta h\right\|_\infty
 \leq\kappa\sup_{0\leq s\leq t}
 \|P_s(\Delta h)-\Delta h\|_\infty
 \longrightarrow0\quad(t\downarrow0).
\]
Ostatnia granica wynika z silnej ciągłości półgrupy.
Dla tych danych generator półgrupy na $\mathrm{BUC}$
działa zatem jako $\kappa\Delta$; nie charakteryzujemy
tutaj całej jego dziedziny.
\end{proof}

\section{Wymuszenie: zasada Duhamela}
\label{sec:pde4-duhamel}

\begin{twierdzenie}[Wzór Duhamela]
\label{thm:pde4-duhamel}
Niech $T>0$, $n\geq1$, $\kappa>0$, $g\in C_c^\infty(\R^n)$,
a $f\in C^\infty([0,T]\times\R^n)$ ma tę własność, że
$f$ i wszystkie jej pochodne względem $x$ i $t$ są
ograniczone, a nośniki $f(t,\cdot)$ są zawarte
we wspólnym zbiorze zwartym $K\subset\R^n$
dla $0\leq t\leq T$. Pochodne czasowe na końcach
przedziału rozumiemy jednostronnie.
Wtedy funkcja
\begin{equation}\label{eq:pde4-duhamel}
 u(t,x)=P_tg(x)+\int_0^t(P_{t-s}f(s,\cdot))(x)
 \,\mathrm d\lambda_1(s),\quad 0\leq t\leq T,
\end{equation}
przy konwencji $P_0h=h$, jest klasy $C^{1,2}$
na $[0,T]\times\R^n$, spełnia
$\partial_tu-\kappa\Delta_xu=f$ oraz $u(0,\cdot)=g$.
Równanie zachodzi również w chwili $0$ w sensie
prawostronnej pochodnej czasowej oraz w chwili $T$
w sensie pochodnej lewostronnej. Ponadto
\begin{equation}\label{eq:pde4-duhamel-bound}
 \|u(t,\cdot)\|_\infty\leq\|g\|_\infty+
 \int_0^t\|f(s,\cdot)\|_\infty\,\mathrm d\lambda_1(s).
\end{equation}
Jest to jedyne rozwiązanie ograniczone na
$[0,T]\times\R^n$ w klasie
$C([0,T]\times\R^n)\cap C^{1,2}((0,T]\times\R^n)$
z tymi danymi.
\end{twierdzenie}
\begin{proof}
Oznaczmy $X=\mathrm{BUC}(\R^n)$ z normą supremum
oraz $F_\alpha(s)=D_x^\alpha f(s,\cdot)$ dla
$\alpha\in\mathbb N_0^n$; w szczególności $F_0(s)=f(s,\cdot)$.
Dla każdego $\alpha$ funkcja $F_\alpha(s)$ należy do $X$,
ponieważ jest ograniczona i ma ograniczony gradient.
Z ograniczoności pochodnej mieszanej
$\partial_sD_x^\alpha f$ otrzymujemy
\[
 \|F_\alpha(s)-F_\alpha(s')\|_\infty
 \leq |s-s'|\sup_{[0,T]\times\R^n}
                  |\partial_tD_x^\alpha f|.
\]
Odwzorowania $F_\alpha:[0,T]\to X$ są więc ciągłe.
Analogiczne warunki przestrzenne spełnia każda
pochodna $g$.

Najpierw uzasadnijmy ciągłość funkcji podcałkowych
w normie $X$ również przy $s=t$. Dla $r,r'\geq0$
oraz $h,h'\in X$ kontrakcyjność daje
\[
 \|P_rh-P_{r'}h'\|_\infty
 \leq\|h-h'\|_\infty+\|P_rh'-P_{r'}h'\|_\infty.
\]
Drugi składnik dąży do zera dla $r\to r'$
przy ustalonym $h'$, dzięki silnej ciągłości półgrupy.
Wynika stąd łączna ciągłość odwzorowania
$(r,h)\mapsto P_rh$ w $[0,\infty)\times X$.
W szczególności dla każdego $\alpha$ funkcja
\[
 H_\alpha(t,s)=P_{t-s}F_\alpha(s),
 \qquad (t,s)\in D_T:=\{(t,s):0\leq s\leq t\leq T\},
\]
jest ciągła w normie $X$ na zwartym trójkącie $D_T$.
Na przekątnej $s=t$ ma wartość $F_\alpha(t)$.
W konsekwencji jest jednostajnie ciągła i ograniczona
na $D_T$. Całki względem $s$ istnieją w $X$ jako
całki funkcji ciągłych i zgadzają się z całkami
punktowymi po obliczeniu wartości w dowolnym $x$.

Dla $r>0$ przenosimy pochodne przestrzenne splotu
na dane, jak w dowodzie lematu
\ref{lem:pde-v31-heat-generator}. Ograniczoność
odpowiednich pochodnych i całkowalność jądra
uzasadniają ten rachunek przez zbieżność zdominowaną.
Dla $r=0$ ta sama tożsamość wynika z definicji $P_0$.
Zatem
$D_x^\alpha(P_rf(s,\cdot))=P_rF_\alpha(s)$
oraz $D_x^\alpha P_rg=P_rD^\alpha g$.
Przy różniczkowaniu całki w~\eqref{eq:pde4-duhamel}
ilorazy różnicowe są ograniczone niezależnie
od $s$ przez supremum odpowiedniej pochodnej
następnego rzędu funkcji $f$ na całym pasku,
na mocy twierdzenia o wartości średniej
i kontrakcyjności $P_{t-s}$.
Dla każdego $t$ dostajemy
\[
 D_x^\alpha u(t,\cdot)
 =P_tD^\alpha g+\int_0^t H_\alpha(t,s)\,
                       \mathrm d\lambda_1(s).
\]
Odwzorowanie po prawej stronie jest ciągłe
w normie $X$ względem $t$: dla $t,t'$ porównujemy
całki na $[0,\min(t,t')]$ przez jednostajną ciągłość
$H_\alpha$ i na pozostałym przedziale przez jego
długość $|t-t'|$ razy ograniczenie $H_\alpha$.
W szczególności w chwili $0$ pochodne przestrzenne
mają wartości $D^\alpha g$.

Pozostaje różniczkowanie względem czasu. Oznaczmy
$F=F_0$, $G(s)=\Delta_x f(s,\cdot)$ oraz
$V(t)=\int_0^tP_{t-s}F(s)\,\mathrm d\lambda_1(s)$.
Z lematu~\ref{lem:pde-v31-heat-generator}
i własności półgrupy wynika dla $0\leq r_1\leq r_2$
oraz $0\leq s\leq T$ tożsamość w $X$
\[
 P_{r_2}F(s)-P_{r_1}F(s)
 =\kappa\int_{r_1}^{r_2}P_qG(s)\,\mathrm d\lambda_1(q).
\]
Istotnie, stosujemy lemat dla $h=F(s)$,
odejmując tożsamości dla $r_2$ i $r_1$;
przy $r_1=0$ korzystamy z $P_0F(s)=F(s)$.
Nie pojawia się więc osobliwość przy $r_1=0$.

Dla $0\leq t<T$ i $0<\tau\leq T-t$ mamy
\[
 \begin{aligned}
 \frac{V(t+\tau)-V(t)}{\tau}
 &=\frac{\kappa}{\tau}\int_0^t
       \int_{t-s}^{t+\tau-s}P_qG(s)\,
       \mathrm d\lambda_1(q)\,\mathrm d\lambda_1(s)\\
 &\quad+\frac1{\tau}\int_t^{t+\tau}P_{t+\tau-s}F(s)\,
                                  \mathrm d\lambda_1(s).
 \end{aligned}
\]
Odwzorowanie $(q,s)\mapsto P_qG(s)$ jest
ciągłe w normie $X$ na zwartym prostokącie
$[0,T]\times[0,T]$, a zatem jednostajnie ciągłe.
Moduł w normie różnicy między pierwszym składnikiem
a $\kappa\int_0^tP_{t-s}G(s)\,\mathrm d\lambda_1(s)$
jest ograniczony przez $\kappa t$ razy supremum
$\|P_qG(s)-P_{t-s}G(s)\|_\infty$ dla
$0\leq s\leq t$ i $|q-(t-s)|\leq\tau$,
$q\in[0,T]$. Oszacowanie to dąży do zera.
Drugi składnik dąży w $X$ do $F(t)$,
ponieważ dla $s\in[t,t+\tau]$ mamy
$0\leq t+\tau-s\leq\tau$ oraz $|s-t|\leq\tau$.
Łączna ciągłość w $(r,s)$ uzasadnia tę granicę
również dla $t=0$.

Dla $0<t\leq T$ i $0<\tau\leq t$ analogicznie
rozpisujemy iloraz lewostronny:
\[
 \begin{aligned}
 \frac{V(t)-V(t-\tau)}{\tau}
 &=\frac{\kappa}{\tau}\int_0^{t-\tau}
      \int_{t-\tau-s}^{t-s}P_qG(s)\,
       \mathrm d\lambda_1(q)\,\mathrm d\lambda_1(s)\\
 &\quad+\frac1{\tau}\int_{t-\tau}^tP_{t-s}F(s)\,
                                   \mathrm d\lambda_1(s).
 \end{aligned}
\]
Pierwszy składnik ma tę samą granicę: oprócz
jednostajnej ciągłości używamy oszacowania brakującej
całki na $[t-\tau,t]$ przez
$\kappa\tau\sup_{s\in[0,T]}\|G(s)\|_\infty$.
Drugi składnik dąży do $F(t)$.
Otrzymujemy dla wszystkich $t\in[0,T]$
(z odpowiednimi pochodnymi jednostronnymi na końcach)
\[
 V'(t)=F(t)+\kappa\int_0^tP_{t-s}G(s)\,
                     \mathrm d\lambda_1(s)
      =F(t)+\kappa\Delta_xV(t).
\]
Pochodna ta jest ciągła w normie $X$ na $[0,T]$
dzięki wcześniej wykazanej ciągłości całek.
W szczególności $V'(0)=F(0)$.

Lemat~\ref{lem:pde-v31-heat-generator} daje również
$\frac{\mathrm d}{\mathrm dt}P_tg=\kappa P_t\Delta g$
w normie $X$ (dla $t=0$ prawostronnie).
Sumując tożsamości, dostajemy
$\partial_tu-\kappa\Delta_xu=f$ i $u(0,\cdot)=g$,
a w chwili początkowej
$\partial_tu(0,\cdot)=\kappa\Delta g+F(0)$.
Ciągłość w normie $X$ funkcji $u$, jej pochodnej
czasowej oraz pochodnych przestrzennych do rzędu
drugiego daje ich ciągłość łączną w $(t,x)$,
przez tę samą nierówność trójkąta, której użyto
w dowodzie twierdzenia~\ref{thm:pde4-cauchy}.
Stąd wynika deklarowana regularność.

Oszacowanie~\eqref{eq:pde4-duhamel-bound} wynika
z kontrakcyjności $P_r$ oraz nierówności trójkąta
dla całki w $X$. Prawa strona jest ograniczona
na $[0,T]$. Różnica dowolnych dwóch ograniczonych
rozwiązań z tym samym wymuszeniem i danymi
początkowymi spełnia jednorodne równanie ciepła
z zerowymi danymi; twierdzenie
\ref{thm:pde4-global-uniqueness} daje jej zerowanie.
\end{proof}

\begin{przyklad}[Źródło stałe w czasie]
Niech $T>0$, $n\geq1$, $\kappa>0$,
$f(t,x)=h(x)$ dla $h\in C_c^\infty(\R^n)$ oraz
$g=0$. Wzór~\eqref{eq:pde4-duhamel}, po podstawieniu
$r=t-s$, daje
\[
 u(t,x)=\int_0^tP_rh(x)\,\mathrm d\lambda_1(r).
\]
Ciągłość $r\mapsto P_rh$ w normie supremum daje
$\partial_tu=P_th$, także prawostronnie w $0$.
Przeniesienie laplasjanu pod całkę oraz na dane,
uzasadnione ograniczonością drugich pochodnych $h$,
prowadzi do
\[
 \kappa\Delta_xu=\kappa\int_0^tP_r\Delta h\,
                       \mathrm d\lambda_1(r)=P_th-h
\]
na mocy lematu~\ref{lem:pde-v31-heat-generator}.
Zatem $\partial_tu-\kappa\Delta_xu=h$,
$u(0,\cdot)=0$ i $\partial_tu(0,\cdot)=h$.
Ponadto $\|u(t,\cdot)\|_\infty\leq t\|h\|_\infty$.
\end{przyklad}

\section{Dyfuzja w obszarze ograniczonym: energia i widmo}
\label{sec:pde4-bounded}

Na obszarze ograniczonym wzór splotowy na całej
przestrzeni na ogół nie spełnia zadanych warunków
brzegowych. Warunki Dirichleta określają temperaturę
na brzegu, a warunki Neumanna --- jej pochodną
normalną, związaną ze strumieniem ciepła przez brzeg.
W tej sekcji wyprowadzimy podstawowe wnioski
dla jednorodnego warunku Dirichleta
\emph{przy założeniu istnienia dostatecznie gładkiego
rozwiązania}. Ogólną teorię istnienia oprzemy później
na przestrzeniach Sobolewa i półgrupach operatorów.

\begin{twierdzenie}[Bilans energii]
\label{thm:pde4-energy}
Niech $n\geq1$, $\Omega\subset\R^n$ będzie obszarem
ograniczonym o brzegu klasy $C^1$, $T>0$, $\kappa>0$.
Niech $u\in C^{1,2}([0,T]\times\overline\Omega)$
będzie funkcją rzeczywistą, spełniającą
$\partial_tu-\kappa\Delta_xu=0$ na $(0,T)\times\Omega$
oraz $u=0$ na $[0,T]\times\partial\Omega$.
Przez podaną klasę regularności rozumiemy, że funkcja
i wskazane pochodne przedłużają się ciągle na
domknięty cylinder, z jednostronnymi pochodnymi
czasowymi na końcach przedziału.
Wówczas dla $0\leq t\leq T$
\begin{equation}\label{eq:pde4-energy}
 \frac12\int_\Omega |u(t,x)|^2\,\mathrm d\lambda_n(x)
 +\kappa\int_0^t\!\int_\Omega|\nabla_xu(s,x)|^2
 \,\mathrm d\lambda_n(x)\,\mathrm d\lambda_1(s)
 =\frac12\int_\Omega|u(0,x)|^2\,\mathrm d\lambda_n(x).
\end{equation}
W szczególności $t\mapsto\|u(t,\cdot)\|_{L^2(\Omega)}$
jest nierosnąca. W powyższej klasie rozwiązanie
z zadanymi danymi początkowymi i zerowym warunkiem
Dirichleta jest jednoznaczne.
\end{twierdzenie}
\begin{proof}
Połóżmy $E(t)=\frac12\int_\Omega |u(t,x)|^2
\,\mathrm d\lambda_n(x)$. Na zwartym zbiorze
$[0,T]\times\overline\Omega$ funkcje $u$ i $\partial_tu$
są ciągłe i ograniczone. Ilorazy różnicowe
$\frac12u(t,x)^2$ względem czasu są zdominowane
przez $\sup|u|\sup|\partial_tu|$, a $\lambda_n(\Omega)$
jest skończona. Z twierdzenia o zbieżności zdominowanej
wynika zatem $E\in C^1([0,T])$, przy jednostronnych
pochodnych na końcach, oraz
$E'(t)=\int_\Omega u(t,x)\partial_tu(t,x)
\,\mathrm d\lambda_n(x)$.
Dla $0<t<T$ podstawiamy równanie ciepła
i stosujemy pierwszą tożsamość Greena
\eqref{eq:pde3-green1} dla ustalonego $t$:
\[
 \begin{aligned}
 E'(t)
 &=\kappa\int_\Omega u(t,x)\Delta_x u(t,x)
             \,\mathrm d\lambda_n(x)\\
 &=-\kappa\int_\Omega|\nabla_xu(t,x)|^2\,\mathrm d\lambda_n(x)\\
 &\quad+\kappa\int_{\partial\Omega}
 u(t,\zeta)\partial_\nu u(t,\zeta)\,\mathrm d\sigma(\zeta).
 \end{aligned}
\]
Tutaj $\nu$ jest jednostkową normalną zewnętrzną,
a $\sigma$ miarą powierzchniową.
W wymiarze $n=1$ ograniczony obszar jest przedziałem
$(a,b)$; ten sam rachunek wynika bezpośrednio
z jednowymiarowego całkowania przez części,
z wyrazem brzegowym
$\kappa[u(t,x)\partial_xu(t,x)]_{x=a}^{x=b}$.
W obu przypadkach wyraz brzegowy znika
dzięki warunkowi Dirichleta.

Funkcja $t\mapsto\int_\Omega|\nabla_xu(t,x)|^2
\,\mathrm d\lambda_n(x)$ jest ciągła na $[0,T]$,
co wynika z ciągłości i ograniczoności gradientu
na domkniętym cylindrze. Mamy więc
\[
 E'(t)=-\kappa\|\nabla_xu(t,\cdot)\|_{L^2(\Omega)}^2.
\]
Równość przedłuża się przez ciągłość do końców
przedziału. Całkując względem czasu od $0$ do $t$,
otrzymujemy~\eqref{eq:pde4-energy}.
Ponieważ $E'\leq0$, energia i norma rozwiązania
są nierosnące.

Jeżeli $u_1,u_2$ spełniają założenia twierdzenia
i mają te same dane początkowe, to $w=u_1-u_2$
spełnia te same warunki z zerowymi danymi.
Bilans energii daje $E_w(t)=0$ dla każdego $t$,
ponieważ obie całki po lewej stronie
są nieujemne, a prawa strona znika.
Gdyby $w(t,x_0)\ne0$ w pewnym $x_0\in\Omega$,
ciągłość dawałaby kulę wewnętrzną, na której
$|w(t,\cdot)|$ jest ograniczona z dołu przez
dodatnią stałą. Przeczyłoby to zerowej normie
$L^2(\Omega)$. Stąd $w=0$ w cylindrze, a na jego
brzegu przestrzennym również z warunku Dirichleta.
\end{proof}

\begin{lemat}[Elementarna nierówność Poincarégo w obszarze ograniczonym]
\label{lem:pde4-poincare}
Niech $n\geq1$, $L>0$ i $\Omega\subset(-L,L)^n$
będzie obszarem ograniczonym o brzegu klasy $C^1$.
Jeżeli $v\in C^1(\overline\Omega)$ jest rzeczywista
i $v|_{\partial\Omega}=0$, to
\begin{equation}\label{eq:pde4-poincare}
 \int_\Omega |v(x)|^2\,\mathrm d\lambda_n(x)
 \leq4L^2\int_\Omega|\nabla v(x)|^2\,\mathrm d\lambda_n(x).
\end{equation}
\end{lemat}
\begin{proof}
Przedłużamy $v$ przez zero do domkniętego sześcianu
$\overline Q=[-L,L]^n$, otrzymując funkcję $V$.
Jest ona ciągła: w punktach brzegu $\Omega$
wartości funkcji $v$ dążą do zera, a poza
$\overline\Omega$ przedłużenie jest lokalnie zerowe.
Wprowadzamy również mierzalną funkcję $W$ na $Q$,
równą $\partial_{x_1}v$ w $\Omega$ i zeru poza
$\Omega$. Jest ona ograniczona, gdyż pochodna
$v$ jest ciągła na zwartym zbiorze $\overline\Omega$.

Dla $n\geq2$ ustalmy $x'\in(-L,L)^{n-1}$
i rozważmy otwarty podzbiór prostej
$I_{x'}=\{s\in(-L,L):(s,x')\in\Omega\}$.
Jest on sumą co najwyżej przeliczalnej rodziny
parami rozłącznych przedziałów otwartych $(a_j,b_j)$.
Każdy ich koniec odpowiada punktowi $\partial\Omega$,
również gdy jest równy $-L$ lub $L$:
jest granicą punktów $\Omega$, ale nie należy
do $\Omega$. Na każdej takiej składowej funkcja
$s\mapsto v(s,x')$ jest klasy $C^1$, a jej wartości
na obu końcach są zerowe. Z twierdzenia podstawowego
rachunku całkowego, najpierw na podprzedziałach
zwartych w $(a_j,b_j)$, a następnie przez
ciągłość na końcach i ograniczoność pochodnej,
wynika
\[
 \int_{a_j}^{b_j}W(s,x')\,\mathrm d\lambda_1(s)=0,
 \qquad
 \int_{a_j}^{q}W(s,x')\,\mathrm d\lambda_1(s)=V(q,x')
 \quad(a_j<q<b_j).
\]
Dla dowolnego $q\in[-L,L]$ możemy rozłożyć całkę
po $[-L,q]$ na przecięcia z tymi składowymi,
ponieważ $W=0$ poza $I_{x'}$, a całka jest
bezwzględnie zbieżna. Wszystkie całe składowe
dają wkład zerowy; co najwyżej jedna składowa
zawiera $q$ i daje wkład $V(q,x')$.
Jeśli $q\notin I_{x'}$, wszystkie wkłady są zerowe,
zgodnie z definicją $V$. Otrzymujemy w każdym przypadku
\[
 V(q,x')=\int_{-L}^{q}W(s,x')\,\mathrm d\lambda_1(s).
\]
W szczególności funkcja $q\mapsto V(q,x')$ jest
absolutnie ciągła na $[-L,L]$ i ma pochodną
równą $W(\cdot,x')$ prawie wszędzie.
Argument nie wymaga, by przekrój miał skończenie
wiele składowych.

Z nierówności Cauchy'ego--Schwarza dostajemy
dla każdego $q\in[-L,L]$
\[
 |V(q,x')|^2\leq 2L
 \int_{-L}^L|W(s,x')|^2\,\mathrm d\lambda_1(s).
\]
Całkując względem $q$ na przedziale długości $2L$
i następnie względem $x'$ na $(-L,L)^{n-1}$,
na mocy twierdzenia Tonellego otrzymujemy
\[
 \int_Q |V|^2\,\mathrm d\lambda_n
 \leq4L^2\int_Q|W|^2\,\mathrm d\lambda_n.
\]
Dla $n=1$ stosujemy bezpośrednio ten sam argument
na przedziale $[-L,L]$, bez zmiennej $x'$.
Z definicji przedłużeń całki po $Q$ są równe
odpowiednim całkom po $\Omega$, a
$|\partial_{x_1}v|^2\leq|\nabla v|^2$.
To daje~\eqref{eq:pde4-poincare}.
\end{proof}

\begin{wniosek}[Wykładniczy zanik energii przy warunku Dirichleta]
\label{cor:pde4-exponential}
Niech $n\geq1$, $T,L,\kappa>0$ i niech
$\Omega\subset(-L,L)^n$ będzie obszarem ograniczonym
o brzegu klasy $C^1$.
Niech $u\in C^{1,2}([0,T]\times\overline\Omega)$
będzie rzeczywistym rozwiązaniem równania
$\partial_tu-\kappa\Delta_xu=0$ na $(0,T)\times\Omega$
z $u=0$ na $[0,T]\times\partial\Omega$.
Wówczas
\begin{equation}\label{eq:pde4-exponential}
 \|u(t,\cdot)\|_{L^2(\Omega)}\leq
 \exp\!\left(-\frac{\kappa t}{4L^2}\right)
 \|u(0,\cdot)\|_{L^2(\Omega)}\qquad(0\leq t\leq T).
\end{equation}
\end{wniosek}
\begin{proof}
Niech $E(t)=\frac12\|u(t,\cdot)\|_{L^2(\Omega)}^2$.
Dla każdego ustalonego $t$ funkcja $u(t,\cdot)$
spełnia założenia lematu~\ref{lem:pde4-poincare}.
Z bilansu energii i~\eqref{eq:pde4-poincare} wynika
\[
 E'(t)=-\kappa\|\nabla_xu(t,\cdot)\|_{L^2(\Omega)}^2
 \leq-\frac{\kappa}{4L^2}\|u(t,\cdot)\|_{L^2(\Omega)}^2
 =-\frac{\kappa}{2L^2}E(t).
\]
Wobec tego pochodna funkcji
$t\mapsto e^{\kappa t/(2L^2)}E(t)$ jest niedodatnia.
Porównując jej wartości w chwilach $0$ i $t$,
otrzymujemy $E(t)\leq e^{-\kappa t/(2L^2)}E(0)$.
Ponieważ $E\geq0$, mnożymy nierówność przez $2$
i pierwiastkujemy, dostając~\eqref{eq:pde4-exponential}.
Zanik wykładniczy dotyczy zarówno energii, jak
i normy $L^2$, ale wykładnik dla normy jest o połowę
mniejszy niż dla energii. Nie twierdzimy,
że stała otrzymana z elementarnej nierówności
Poincarégo jest optymalna.
\end{proof}

\begin{przyklad}[Mody własne na odcinku]
Niech $\Omega=(0,L)$, gdzie $L>0$, oraz $\kappa>0$.
Dla każdego całkowitego $m\geq1$ określamy
$\phi_m(x)=\sin(m\pi x/L)$ i $\lambda_m=(m\pi/L)^2$.
Mamy $\phi_m''=-\lambda_m\phi_m$ oraz
$\phi_m(0)=\phi_m(L)=0$.
Funkcja
\[
 u_m(t,x)=e^{-\kappa\lambda_m t}\phi_m(x)
\]
spełnia $\partial_tu_m=-\kappa\lambda_m u_m$
i $\partial_x^2u_m=-\lambda_m u_m$, zatem
rozwiązuje równanie ciepła, znika na brzegu
i przyjmuje w chwili $0$ dane $\phi_m$.
Ponieważ $\sin^2\theta=(1-\cos(2\theta))/2$,
mamy $\|\phi_m\|_{L^2(0,L)}^2=L/2$.
Stąd
\[
 \|u_m(t,\cdot)\|_{L^2(0,L)}^2
 =\frac L2 e^{-2\kappa\lambda_m t},\qquad
 \|\partial_xu_m(t,\cdot)\|_{L^2(0,L)}^2
 =\lambda_m\frac L2 e^{-2\kappa\lambda_m t}.
\]
Różniczkując połowę pierwszej równości względem
czasu, otrzymujemy dokładnie bilans energii
z twierdzenia~\ref{thm:pde4-energy}.
W szczególności większe $\lambda_m$ odpowiadają
szybszemu zanikowi amplitudy.

Z liniowości równania dla każdego $N\geq1$
i $c_1,\ldots,c_N\in\R$ funkcja
$\sum_{m=1}^N c_m u_m$ również jest rozwiązaniem,
z danymi początkowymi $\sum_{m=1}^N c_m\phi_m$.
Jego jednoznaczność wynika z bilansu energii.
Ogólne rozwinięcia w szereg funkcji własnych
laplasjanu wymagają ponadto uzasadnienia
zbieżności i możliwości różniczkowania szeregu;
wrócimy do nich po wprowadzeniu teorii Fouriera.
\end{przyklad}

\section{Zachowanie dla dużych czasów i nieodwracalność}
\label{sec:pde4-large-time}

\begin{twierdzenie}[Zanik na całej przestrzeni]
\label{thm:pde4-large-time}
Niech $n\geq1$, $\kappa>0$.
Jeśli $g\in L^1(\R^n)$, to $\|P_tg\|_\infty\to0$
przy $t\to\infty$.
Jeśli $g\in L^2(\R^n)$, to $\|P_tg\|_2\to0$
przy $t\to\infty$.
\end{twierdzenie}
\begin{proof}
Dla $g\in L^1(\R^n)$ pierwsze oszacowanie
w~\eqref{eq:pde4-L1Linf} daje
\[
 0\leq\|P_tg\|_\infty\leq
 (4\pi\kappa t)^{-n/2}\|g\|_1\longrightarrow0,
\]
ponieważ $n\geq1$ i $\kappa>0$.
Jeżeli $g\in L^1(\R^n)\cap L^2(\R^n)$, drugie
oszacowanie w~\eqref{eq:pde4-L1Linf} daje
\[
 0\leq\|P_tg\|_2\leq
 (8\pi\kappa t)^{-n/4}\|g\|_1\longrightarrow0.
\]

Niech teraz $g\in L^2(\R^n)$ będzie dowolna.
Ustalmy $\eps>0$. Z gęstości $C_c^\infty(\R^n)$
w $L^2(\R^n)$ wybieramy $h\in C_c^\infty(\R^n)$
takie, że $\|g-h\|_2<\eps/2$.
Funkcja $h$ należy również do $L^1(\R^n)$,
więc istnieje $T_\eps>0$ takie, że
$\|P_th\|_2<\eps/2$ dla wszystkich $t\geq T_\eps$.
Z liniowości i kontrakcyjności $P_t$ w $L^2$,
wykazanej w twierdzeniu~\ref{thm:pde4-approx-identity},
otrzymujemy dla $t\geq T_\eps$
\[
 \|P_tg\|_2
 \leq\|P_t(g-h)\|_2+\|P_th\|_2
 \leq\|g-h\|_2+\|P_th\|_2<\eps.
\]
Dowolność $\eps$ daje żądaną zbieżność.
Argument wymaga tylko aproksymacji danych w $L^2$,
nie zakłada ich całkowalności w $L^1$.
\end{proof}

Należy odróżnić silny zanik dla każdego ustalonego
$g\in L^2$ od zaniku normy operatorowej półgrupy
na $L^2$. Powyższy dowód nie daje wspólnego tempa
zaniku dla wszystkich danych o normie $1$.
Istotnie, ustalmy niezerową funkcję
$h\in C_c^\infty(\R^n)$ z $\|h\|_2=1$ i dla $R>0$
połóżmy $h_R(x)=R^{-n/2}h(x/R)$.
Zmiana zmiennych $x=Ry$ daje $\|h_R\|_2=1$.
Ze skalowania jądra i podstawienia $y=Rz$ w splocie
wynika
\[
 (P_th_R)(x)=R^{-n/2}(P_{t/R^2}h)(x/R).
\]
Zatem $\|P_th_R\|_2=\|P_{t/R^2}h\|_2\to1$
przy $R\to\infty$ dla każdego ustalonego $t>0$,
na mocy silnej ciągłości w chwili zero.
W połączeniu z kontrakcyjnością daje to
$\|P_t\|_{L^2\to L^2}=1$ dla każdego $t\geq0$.
Na całej przestrzeni nie otrzymujemy więc
jednolitego wykładniczego zaniku analogicznego
do oszacowania~\eqref{eq:pde4-exponential}
na ograniczonym obszarze z warunkiem Dirichleta.

Zanik nie musi zachodzić także dla ogólnych danych
ograniczonych: dla $g\equiv1$ normalizacja jądra
daje $P_tg\equiv1$. Z kolei dla nieujemnego
$g\in L^1(\R^n)$, $g\not=0$ w $L^1$, norma $L^1$
pozostaje stała z własności zachowania masy
\eqref{eq:pde4-mass}. Zanik normy supremum
dla takich danych nie oznacza zaniku całkowitej masy.

\begin{kontrprzyklad}[Wsteczne równanie ciepła jest niestabilne]
\label{ex:pde4-backward}
Niech $\Omega=(0,\pi)$, $\kappa=1$ i $T>0$.
Rozważmy przestrzeń $S$ skończonych kombinacji
liniowych funkcji $\phi_m(x)=\sin(mx)$,
gdzie $m\geq1$ jest liczbą całkowitą.
Z przykładu modów własnych oraz jednoznaczności
w twierdzeniu~\ref{thm:pde4-energy} wynika, że
odwzorowanie ewolucji od chwili $0$ do $T$ na $S$
jest dobrze określone i ma postać
\[
 A_T\left(\sum_{m=1}^N c_m\phi_m\right)
 =\sum_{m=1}^N e^{-m^2T}c_m\phi_m.
\]
Funkcje $\phi_m$ są liniowo niezależne:
z tożsamości
$2\sin(mx)\sin(kx)=\cos((m-k)x)-\cos((m+k)x)$
otrzymujemy dla $m\ne k$
$\int_0^\pi\phi_m(x)\phi_k(x)\,\mathrm d\lambda_1(x)=0$,
a $\|\phi_m\|_{L^2(0,\pi)}^2=\pi/2$.
Ponieważ wszystkie mnożniki $e^{-m^2T}$ są dodatnie,
$A_T:S\to S$ jest liniową bijekcją.

Dla danych początkowych $\phi_m$ mamy
$u_m(t,x)=e^{-m^2t}\sin(mx)$.
Stąd
\[
 \|\phi_m\|_\infty=1,\qquad
 \|A_T\phi_m\|_\infty=e^{-m^2T}\longrightarrow0.
\]
Jeżeli $v_m=A_T\phi_m$, to $v_m\to0$ w normie
supremum, natomiast $A_T^{-1}v_m=\phi_m$
nie dąży do zera w tej normie.
Odwrotne odwzorowanie $A_T^{-1}:S\to S$
jest więc nieciągłe w zerze względem normy supremum.

Również w normie $L^2(0,\pi)$ mamy
\[
 \|v_m\|_2=e^{-m^2T}\sqrt{\pi/2}\longrightarrow0,
 \qquad \|A_T^{-1}v_m\|_2=\sqrt{\pi/2}.
\]
W obu normach iloraz
$\|A_T^{-1}v_m\|/\|v_m\|=e^{m^2T}$ jest
nieograniczony. Nawet ograniczenie do dziedziny $S$
nie pozwala więc uzyskać ciągłego odwracania.
Każde odwrotne odwzorowanie na większej klasie
danych końcowych, zgodne z odwracaniem tych modów,
musiałoby ograniczać się do $A_T^{-1}$ na $S$,
a zatem również nie mogłoby być ciągłe w zerze
w żadnej z tych dwóch norm.

Jest to przykład niestabilności w sensie Hadamarda:
nawet w klasie, w której istnienie i jednoznaczność
odwrócenia są zapewnione, ciągła zależność danych
początkowych od końcowych zawodzi.
Nie wynika stąd, że dwa różne dane z $S$
prowadzą do tego samego profilu końcowego.
Termin nieodwracalność oznacza tu brak stabilnego
odwracania w rozpatrywanych normach, nie brak
algebraicznej odwrotności na obrazie ewolucji.
\end{kontrprzyklad}

\section{Nieoczekiwana niejednoznaczność: przykład Tychonoffa}
\label{sec:pde4-tychonoff}

Powyższa jednoznaczność wymagała ograniczoności
rozwiązania w całym cylindrze czasoprzestrzennym.
Pokażemy teraz, że nie jest to wyłącznie techniczny
warunek dowodu.

\begin{twierdzenie}[Istnienie niezerowego rozwiązania o zerowych danych]
\label{thm:pde4-tychonoff}
Istnieje niezerowa funkcja
$u\in C^\infty([0,\infty)\times\R)$
spełniająca $u_t=u_{xx}$ dla $t>0$, $x\in\R$,
przy czym $u(0,x)=0$ dla wszystkich $x\in\R$.
Funkcja ta nie może być ograniczona na żadnym
cylindrze $[0,T]\times\R$ z $T>0$.
\end{twierdzenie}
\begin{proof}
Niech $h(t)=e^{-1/t^2}$ dla $t>0$ i $h(t)=0$
dla $t\leq0$. Na dodatniej półosi funkcja $h$
jest klasy $C^\infty$. Najpierw definiujemy formalnie
szereg, a następnie uzasadnimy jego zbieżność
i gładkość sumy także przy $t=0$. Kładziemy
\begin{equation}\label{eq:pde4-tychonoff}
 u(t,x)=\sum_{k=0}^\infty\frac{h^{(k)}(t)}{(2k)!}
 x^{2k}\quad(t>0),\qquad u(0,x)=0.
\end{equation}
Uzasadnijmy najpierw oszacowanie pochodnych $h$.
Dla $t>0$ i $z=a+ib$ spełniającego $|z-t|\leq t/4$
mamy $a\geq3t/4$, $|b|\leq t/4$ oraz $|z|\leq5t/4$.
Wobec tego
\[
 \operatorname{Re}\frac1{z^2}
 =\frac{a^2-b^2}{|z|^4}
 \geq\frac{t^2/2}{(5t/4)^4}
 =\frac{128}{625t^2}.
\]
Możemy przyjąć $c=128/625>0$, niezależne od $t$.
Funkcja $z\mapsto e^{-1/z^2}$ jest holomorficzna
na otoczeniu domkniętego dysku o środku $t$
i promieniu $t/4$, gdyż dysk nie zawiera zera.
Na jego brzegu moduł tej funkcji nie przekracza
$e^{-c/t^2}$. Oszacowanie Cauchy'ego na tym okręgu
daje dla każdego całkowitego $k\geq0$
\begin{equation}\label{eq:pde4-tychonoff-est}
 |h^{(k)}(t)|\leq k!\left(\frac4t\right)^k
 e^{-c/t^2},\qquad k\geq0.
\end{equation}
Ponieważ $\binom{2k}{k}\geq1$, mamy $(2k)!\geq(k!)^2$.
Dla każdego $A>0$ i $|x|\leq A$
otrzymujemy
\[
 |u(t,x)|\leq e^{-c/t^2}
 \sum_{k=0}^\infty\frac{(4A^2/t)^k}{k!}
 =\exp\!\left(-\frac c{t^2}+\frac{4A^2}t\right).
\]
Majoranta wykładnicza jest skończona dla każdego
$t>0$, więc szereg definiujący $u$ jest bezwzględnie
zbieżny. Oszacowanie dowodzi także jednostajnej
zbieżności $u(t,\cdot)$ do zera na $[-A,A]$
przy $t\downarrow0$.

Uzasadnijmy dokładniej oszacowanie wszystkich pochodnych.
Dla ustalonych $p,q\in\mathbb N_0$ oraz $2k\geq q$ mamy
\[
 \partial_t^p\partial_x^q
 \left(\frac{h^{(k)}(t)x^{2k}}{(2k)!}\right)
 =\frac{h^{(k+p)}(t)x^{2k-q}}{(2k-q)!}.
\]
Jeżeli $2k<q$, pochodna jest identycznie zerowa.
Wybierzmy $B=\max\{1,A\}$ dla $|x|\leq A$.
Oszacowanie~\eqref{eq:pde4-tychonoff-est} daje
\[
 \left|\frac{h^{(k+p)}(t)x^{2k-q}}{(2k-q)!}\right|
 \le e^{-c/t^2}\left(\frac4t\right)^p
 \frac{(k+p)!}{(2k-q)!}
 \left(\frac{4B^2}{t}\right)^k.
\]
Dla $k\geq1$ i $2k\geq q$ mamy
\[
 \frac{(k+p)!}{(2k-q)!}
 =\frac{(k+p)!}{k!}\frac{k!}{(2k)!}
   \frac{(2k)!}{(2k-q)!}
 \leq\frac{(k+p)^p(2k)^q}{k!}
 \leq C_{p,q}\frac{(1+k)^{p+q}}{k!},
\]
gdzie pusty iloczyn ma wartość $1$.
Dla $k=0$ dopuszczalny jest tylko przypadek $q=0$;
zwiększamy wtedy stałą tak, by obejmowała $p!$.

Dla $m\in\mathbb N_0$ i $z\geq0$ mamy
$\sum_{k\geq0}(1+k)^m z^k/k!=Q_m(z)e^z$,
gdzie $Q_0=1$ oraz
$Q_{m+1}=(1+z)Q_m+zQ_m'$.
Istotnie, stosujemy operator $1+z\partial_z$
do szeregu, co jest dozwolone, gdyż ma on
nieskończony promień zbieżności (z kryterium ilorazowego).
Rekurencja pokazuje, że $Q_m$ ma stopień $m$
i nieujemne współczynniki, więc
$Q_m(z)\leq C_m(1+z)^m$ dla $z\geq0$.
Stąd
cały szereg pochodnych jest ograniczony na
$0<t\le1$, $|x|\le A$ przez
\[
 C_{p,q,A}t^{-p}(1+t^{-1})^{p+q}
 \exp\left(-\frac c{t^2}+\frac{4B^2}{t}\right).
\]
Najpierw uzasadnijmy różniczkowanie dla $t>0$.
Na prostokącie $[a,b]\times[-A,A]$, $0<a<b$,
każdy składnik szeregu pochodnych jest ograniczony
przez
\[
 C_{p,q}(4/a)^p
 \frac{(1+k)^{p+q}}{k!}(4B^2/a)^k .
\]
Szereg tych majorant jest zbieżny z kryterium
ilorazowego. Kryterium Weierstrassa daje więc
jednostajną zbieżność każdego szeregu pochodnych
na takim prostokącie. Twierdzenie o różniczkowaniu
szeregów funkcyjnych, stosowane kolejno na mniejszych
otwartych prostokątach, daje $u\in C^\infty$
dla $t>0$ i identyfikuje każdą jej pochodną
$\partial_t^p\partial_x^q u$ z odpowiednim szeregiem.

Sprawdźmy teraz gładkość w chwili zero.
Dla dostatecznie małych $t>0$, konkretnie
$t\leq c/(8B^2)$, mamy
$-c/t^2+4B^2/t\leq-c/(2t^2)$.
Dla każdego $N\geq0$ i ustalonych $p,q,A$
otrzymujemy z wcześniejszego oszacowania
\[
 \sup_{|x|\leq A}|\partial_t^p\partial_x^q u(t,x)|
 =o(t^N)\qquad(t\downarrow0),
\]
gdyż wykładniczy czynnik $e^{-c/(2t^2)}$
zanika szybciej od każdej potęgi $t$.

Przedłużmy $u$ przez zero na $t<0$ i oznaczmy
dla każdego $p,q$ przez $U_{p,q}$ funkcję równą
$\partial_t^p\partial_x^q u$ dla $t>0$ oraz zeru
dla $t\leq0$. Każda $U_{p,q}$ jest ciągła
na $\R^2$ dzięki lokalnie jednostajnemu oszacowaniu
z $N=0$. W punkcie $(0,x)$ iloraz różnicowy
$U_{p,q}$ względem czasu od lewej jest zerowy,
a od prawej dąży do zera na mocy oszacowania
z $N=1$. Względem $x$ pochodna w chwili zero
jest zerowa, ponieważ $U_{p,q}(0,\cdot)\equiv0$.
Stąd na całej $\R^2$
\[
 \partial_tU_{p,q}=U_{p+1,q},\qquad
 \partial_xU_{p,q}=U_{p,q+1}.
\]
Indukcyjnie wynika, że $U_{0,0}$ jest klasy
$C^\infty(\R^2)$ i wszystkie jej pochodne
znikają na $\{0\}\times\R$.
W szczególności jej ograniczenie do $t\geq0$
ma wymaganą gładkość.

Dla $t>0$ różniczkowanie wyraz po wyrazie daje
\[
 \begin{aligned}
 u_t(t,x)&=\sum_{k=0}^\infty
     \frac{h^{(k+1)}(t)}{(2k)!}x^{2k},\\
 u_{xx}(t,x)&=\sum_{k=1}^\infty
     \frac{h^{(k)}(t)}{(2k-2)!}x^{2k-2}\\
 &=\sum_{j=0}^\infty\frac{h^{(j+1)}(t)}{(2j)!}x^{2j}.
 \end{aligned}
\]
Zatem $u_t=u_{xx}$, również przy $t=0$,
gdzie obie pochodne znikają.
Wreszcie $u(t,0)=h(t)>0$ dla każdego $t>0$,
więc nie jest to rozwiązanie zerowe.
Gdyby było ograniczone na pewnym cylindrze $[0,T]\times\R$,
twierdzenie~\ref{thm:pde4-global-uniqueness}
wymusiłoby $u=0$, co jest sprzeczne z $u(t,0)>0$.
\end{proof}

\begin{uwaga}[Co ten kontrprzykład rzeczywiście oznacza?]
Twierdzenie~\ref{thm:pde4-global-uniqueness}
zapewnia jednoznaczność w klasie rozwiązań
ograniczonych na każdym pasku $[0,T]\times\R^n$.
Istnieją również twierdzenia o jednoznaczności
przy innych warunkach wzrostu, których tutaj
nie formułujemy. Przykład Tychonoffa pokazuje,
że sama gładkość na domkniętym pasku i zerowe
dane początkowe nie wystarczają.
Nie wynika z niego, że każde rozwiązanie
nieograniczone prowadzi do niejednoznaczności.
Ponieważ rozwiązanie i wszystkie jego pochodne
znikają w chwili zero, a $u(t,0)>0$ dla $t>0$,
nie jest ono analityczne względem czasu
w otoczeniu chwili początkowej.
Nie przeczy to twierdzeniu o analityczności
przestrzennej splotu z jądrem Gaussa.
\end{uwaga}

\section{Podsumowanie i związek z kolejnym rozdziałem}
\label{sec:pde4-summary}

Kluczowe mechanizmy równania ciepła są widoczne
w~\eqref{eq:pde4-kernel}: dodatniość jądra daje
zachowanie porządku dla rozwiązań splotowych
i natychmiastową propagację niezerowych danych
nieujemnych o zwartym nośniku. Normalizacja
zapewnia zachowanie całki dla danych z $L^1$,
a skalowanie Gaussa wyjaśnia oszacowania zaniku.
Półgrupa jest silnie ciągła na $L^p$ dla $p<\infty$
($p\geq1$) oraz na $\mathrm{BUC}$, lecz nie na
całym $L^\infty$.

Jawny wzór rozwiązuje jednorodne zagadnienie
Cauchy'ego z danymi w $\mathrm{BUC}$, a wzór
Duhamela obejmuje niejednorodne równanie przy
podanych założeniach gładkości danych i wymuszenia.
Jednoznaczność klasycznych rozwiązań na całej
przestrzeni wymaga określenia klasy rozwiązań:
dla funkcji ograniczonych na każdym pasku
czasowym zapewnia ją zasada maksimum, natomiast bez
kontroli wzrostu przykład Tychonoffa daje
niejednoznaczność.

Dla obszarów ograniczonych wyprowadziliśmy bilans
energii i wykładniczy zanik normy przy jednorodnym
warunku Dirichleta, zakładając istnienie odpowiednio
gładkiego rozwiązania. Nie stanowi to jeszcze
ogólnego twierdzenia o istnieniu.
Na całej przestrzeni silny zanik w $L^2$ nie oznacza
zaniku normy operatorowej półgrupy. Wygładzanie
przestrzenne oraz tłumienie wysokich częstości
prowadzą do niestabilności odtwarzania danych
początkowych; nie oznacza ona braku algebraicznej
odwrotności na przestrzeni skończonych sum modów.
W następnym rozdziale porównamy te fakty z równaniem
falowym, gdzie propagacja ma skończoną prędkość,
a klasyczna energia ma inny charakter.


% Rozdział 5. Równanie falowe.
